\section{灵敏度分析与误差分析}

为检验模型结果对关键参数与输入误差的依赖程度，本章从灵敏度分析与误差分析两方面进行检验：
前者考察关键参数变化时最优方案与购电费用的变化幅度，后者考察预测值与实际值的偏差及其经济后果。

\subsection{灵敏度分析}

\subsubsection{问题一：日循环边界与充放电效率}
问题一的灵敏度检验结果如表~\ref{tab:sens-q1} 所示。

\begin{table}[!htbp]
\caption{问题一灵敏度分析（全天购电费用）}\label{tab:sens-q1}
\centering
\small
\begin{tabular}{lcc}
\toprule
情形 & 全天购电费/元 & 与基准差 \\
\midrule
基准（$S(0)=S(144)=6000$\,kWh） & 35126.95 & --- \\
首末储电量自由（$\to8550$\,kWh，附录 A） & 35118.60 & $-$8.35（$-$0.024\%） \\
往返效率 90\%（$\eta_c=\eta_d=\sqrt{0.9}$） & 33793.78 & $-$3.8\% \\
无储能设备 & 48052.05 & $+$36.8\% \\
\bottomrule
\end{tabular}
\end{table}

日循环边界的两种解读对费用影响仅 0.024\%，结论稳健；
充放电效率由 81\% 往返提升至 90\% 往返可使费用进一步下降 3.8\%，
说明效率参数是储能经济性的重要来源；无储能时费用上升 36.8\%，再次确认储能的核心价值。
约束取等情况：储能上限 10800\,kWh 在 14:30--18:00 与 5:40--5:50 取等、下限 1200\,kWh 在 20:50--22:00 取等，
充电功率 5000\,kW 在 19 个时段满功率取等——对偶意义下扩大储能容量与充电通道具有正的边际收益。
此外，无爬坡成本时 LP 最优解呈“bang-bang”特征（购电曲线相邻区间最大跳变 4994\,kW、
$>$500\,kW 跳变 22 次）；对购电功率加微小爬坡惩罚 $\varepsilon\sum_t|p(t)-p(t-1)|$（仍为 LP），
$\varepsilon=0.005$ 元/kW 时费用仅增加 67.12 元（$+$0.19\%），跳变降至 9 次，
实际申报时可采用该平滑化方案。

\subsubsection{预测方法与余量参数}
模型中不确定性最强、对结果影响最大的是预测方法与预测余量系数。
表~\ref{tab:sens-q2} 给出问题二在不同预测方法与余量组合下的全年总费用与紧急购电费用。

\begin{table}[!htbp]
\caption{不同预测方法与余量组合下的全年运行费用（问题二，报告期 334 天）}\label{tab:sens-q2}
\centering
\small
\setlength{\tabcolsep}{5pt}
\begin{tabular}{llcc}
\toprule
预测方法 & $(\gamma,\delta)$ & 总费用/万元 & 紧急购电费/万元 \\
\midrule
持续性预测（当日 $=$ 昨日） & $(1.00,1.00)$ & 2244.8 & 1021.1 \\
持续性预测 & $(1.05,1.00)$ & 2139.4 & 793.6 \\
上周同星期（w1） & $(1.00,1.00)$ & 1602.4 & 378.7 \\
\textbf{上周同星期（w1，本文采用）} & $\bm{(1.05,0.98)}$ & \textbf{1454.7} & \textbf{87.0} \\
近 4 个同星期均值 & $(1.05,1.00)$ & 1574.7 & 229.7 \\
7 日滑动平均 & $(1.08,1.00)$ & 2059.0 & 698.6 \\
\bottomrule
\end{tabular}
\end{table}

由表~\ref{tab:sens-q2} 可知，预测方法的选取对费用的影响远大于余量系数：
持续性预测因周末前后的负载状态切换而严重失效（2244.8 万元）；
w1 预测不加余量时为 1602.4 万元，加入余量后降至 1454.7 万元（降 147.7 万元），
说明\textbf{余量策略是应对 5 倍紧急购电惩罚的关键}。
w1 预测下 $\gamma$ 在 $[1.03,1.07]$ 内构成平坦平台（偏移 $\pm0.02$ 时总费用变化不足 1\%）；
将报告期拆分为 2--8 月与 9--12 月两段分别寻优（时间序列交叉验证\cite{fpp3}），
最优参数与全年结果一致，说明参数并非对全年数据过拟合。

\subsubsection{策略参数的稳健性}
\textbf{（1）储能配置的价值。} 同一预测与余量下取消储能设备，
问题二全年总费用由 1454.7 万元升至 1893.9 万元，储能年价值 439.2 万元（$-$23.2\%）。

\textbf{（2）当日残差更新权重与日末备用。} 问题三中 $\lambda=0$ 时总费用 1376.3 万元，
最优 $(0.1,0.3)$ 处 1368.8 万元，增大至 $(0.6,0.8)$ 回升至 1414.4 万元（表~\ref{tab:q3-lambda}），
最优点附近平缓；日末最低备用在 2000$\sim$7000\,kWh 范围内总费用变化平缓，
不改变“标定 $+$ 双向调整优于无调整”的结论，取 4000\,kWh 主要出于备用稳健性考量。

\textbf{（3）波动电价下的参数复核。} 引入波动电价（问题四）后重新扫描余量参数，
结果如表~\ref{tab:sens-q4} 所示：Q4-2 最优值仍在 $\gamma=1.05$ 附近，继续沿用 $(1.05,0.98)$；
Q4-3 负载侧最优值仍为 $\gamma=1.02$，光伏侧最优余量由 $\delta=0.98$ 变为 $\delta=1.00$——
波动电价下光伏边际价值下降，过度下压光伏预测反而抬高购电成本。
电价波动改变了最优光伏余量，但未改变模型结构与整体结论。

\begin{table}[!htbp]
\caption{波动电价下余量参数复核（全年总费用，单位：万元）}\label{tab:sens-q4}
\centering
\small
\begin{tabular}{lccc}
\toprule
Q4-2：$\gamma$（$\delta=0.98$） & 1.03 & \textbf{1.05} & 1.07 \\
\midrule
全年总费用 & 1530.2 & \textbf{1526.4} & 1548.6 \\
\midrule
Q4-3：$\delta$（$\gamma=1.02$） & 0.96 & 0.98 & \textbf{1.00} \\
\midrule
全年总费用 & 1440.3 & 1433.6 & \textbf{1431.6} \\
\bottomrule
\end{tabular}
\end{table}

\subsection{误差分析}

\subsubsection{光伏预报误差}
对附件 3 整点预报与实际出力的对照（图~\ref{fig:q3-forecast}）表明，
白天时段预报 MAE 最高达 1374.1\,kW，并呈现稳定的\textbf{系统性偏差}：
上午偏低（8:00 约 $-1050$\,kW）、午后偏高（16:00 约 $+988\sim+1067$\,kW），
预报峰值时刻与实际峰值时刻错位。经 MOS 标定后，14:00 的估计 MAE 由原始预报的 793.4\,kW
降至 373.0\,kW，也优于单独使用上周同星期实际值的 456.8\,kW，
说明\textbf{标定有效削弱了系统性偏差}。由于预报误差不可能完全消除，
本文并未追求更高的点预测精度，而是通过安全余量与日内滚动调整吸收剩余误差。

\subsubsection{预测误差的经济代价}
以“完美预测”为理论下界：在与最终策略完全相同的框架下（含储能、日循环约束与跨日滚动），
仅将预测值替换为当日实际值并将余量置 1，重新求解全年调度。
完美预测下全年总费用为 1224.5 万元，本文问题二策略为 1454.7 万元，
二者相差 230.2 万元，即\textbf{残余预测误差的代价约为总费用的 15.8\%}；
问题三引入标定与滚动调整后，该差距缩小至 144.3 万元（1368.8 万元对 1224.5 万元），
说明滚动调整机制确实降低了预测误差的经济损失。
进一步考察误差代价的分布：问题二紧急购电的平均结算电价为 1.094 元/kWh，
约为全天均价的 1.4 倍——供电缺口较多落在傍晚高电价时段，5 倍惩罚在最贵时段被放大，
这正是问题三滚动调整收益的主要来源。

\subsubsection{余量策略的固有代价与数值校验}
日前计划施加安全余量的代价是部分已付费电量最终无法使用：
问题二全年弃电量约 326 万\,kWh（占计划购电量 14.6\%），
其中“已付费却被迫丢弃”的电量下限约 81 万\,kWh，按均价估算浪费不低于 60 万元（约占总费用 4\%）。
这是“日前承诺 $+$ 不确定性对冲”的必然成本，与所避免的约 790 万元紧急购电费用相比净收益显著。

最后，对求解结果作全链路数值校验：（1）各时段功率平衡与储能状态递推均成立；
（2）储能轨迹全程位于 $[1200,10800]$\,kWh，充放电功率不超过 5000\,kW；
（3）跨日储电量按实际末态连续传递，无跳变；
（4）各日费用按“计划购电费 $+$ 调整费用 $+$ 紧急购电费”分解后与总费用一致；
（5）调整模型的解可由重新求解线性规划复现。
上述校验在 334 天滚动仿真中均未发现违规；全部数值结果经 Python 与 MATLAB 双平台独立求解逐位一致，
说明求解结果自洽且可复现。
